\section{Welfare and sale terms}\label{sec:welfare}

\subsection{Access to prices}\label{sec:access}

Access to prices improves acquisition outcomes at a fixed level of competition. At the strong incumbent
\(F_1\), I compare the equilibrium of Proposition~\ref{prop:ccc} with that of an economy in which the
rival does not see the price. Under the hypotheses of Proposition~\ref{prop:ccc}, both economies have unique full orders, and price access
raises expected target proceeds and acquisition surplus net of preparation costs
(Proposition~\ref{prop:access}).
Acquisition surplus excludes transfers between bidders, shareholders and traders, and
Appendix~\ref{app:welfare} states each result on welfare and sale terms in full.

The two economies generate the same orders. The trading bound of Section~\ref{sec:informed} forces full orders even
when the rival cannot use the price and entry stays at \(\rho\), so trading costs coincide. Price access
changes only which high-cost rivals prepare: it adds preparation after a purchase. Each added preparation
adds its gross profit minus its cost to net surplus, together with the value of any sale that the reserve
would otherwise block; a rival prepares only when its expected gross profit covers its cost, so no
added preparation lowers net surplus.

In the example at \(r_1\), expected proceeds rise from \(\reg{p2_base_hidden_proceeds}\) to
\(\reg{p2_base_r1_proceeds}\), a gain of \(\reg{p2_access_proceeds_gain}\), and net surplus rises by
\(\reg{p2_access_surplus_gain}\). With research at the cap (Section~\ref{sec:research}), the
gains are the same.

Target shareholders gain because a shareholder who holds through the sale receives \(V_T\), not the interim price.
The informed trader's expected gross profit, \(\reg{p2_access_gross_transfer}\) at \(r_1\), is a transfer from
liquidity traders. If one order unit is a given number of shares per share of float, initial shareholders gain in
total relative to an economy without the informed trader whenever that number is below \(\reg{p2_access_break_even}\), even if
they bear the whole transfer.

Inside the recruitment region of Proposition~\ref{prop:ccc}, where high-cost preparation does not pay at the prior,
price access adds preparation. When gross profit at the prior exceeds \(c_H\) and gross profit after a sale falls below it, an informative price
deters high-cost preparation after a sale: in the example at \(r_1\) with \(c_H=\reg{access_out_cH}\), price access changes proceeds by
\(\reg{p2_outside_proceeds_gain}\) and net surplus by \(\reg{p2_outside_surplus_gain}\)
(Appendix~\ref{app:welfare}).

The two economies differ in the information in the price and not in its level. A deterministic external dividend of
\(\reg{p2_access_proceeds_gain}\) in the hidden-price economy equalizes the mean prices, and it leaves the orders
and entry unchanged.

\subsection{The payment rule}\label{sec:payment}

The effect of competition on the target-payoff spread depends on the sale payment.
Proposition~\ref{prop:payment} isolates this dependence at the payment stage, in a
verifiable-value institution with zero reserve, point-mass classes and an incumbent distribution on \([0,\bar r]\)
with \(0<\ell<\bar r<h\). After diligence, a sale to the next bidder at its value is an
enforceable fallback, and the seller and the winner divide the surplus above the fallback by Nash bargaining, with
seller weight \(\omega\in[0,1)\). A stronger incumbent never raises the rival's profit, and it moves the
spread only in the direction of the sign of \(1-2\omega\); the move is strict when \(\E[(R-\ell)_+]\) changes
strictly and \(1-2\omega\ne0\), and when \(1-2\omega=0\) the spread equals \((h-\ell)/2\) at every incumbent
distribution.

To see why, take \(R>\ell\). A high-class rival wins and pays \((1-\omega)R+\omega h\), whose sensitivity to
\(R\) is \(1-\omega\). A low-class rival loses, and the incumbent pays \((1-\omega)\ell+\omega R\), whose
sensitivity is \(\omega\). Competition widens the difference exactly when the first sensitivity exceeds the second.

In the example, the spread under this rule changes by \(\reg{pay_gap_coef}\,(1-2\omega)\) from \(r_0\) to \(r_1\).
The second-price auction is the case \(\omega=0\): it puts the whole weight on the runner-up's value.
\citet{BulowKlemperer1996} compare an auction with one more bidder and an optimal negotiation by expected revenue;
here the comparison concerns how competition moves the spread.

\subsection{Reserves}\label{sec:reserves}

A reserve affects both the payment required to win and the information in the target's shares. Once the reserve can
exclude bidders, preparation, a sale and the meeting of two admissible bidders become different events:
\emph{preparation} is the event that the rival pays its preparation cost, a \emph{sale} occurs when at
least one bidder submits an admissible bid, and \emph{two admissible bidders} meet when both the incumbent
and a prepared rival meet the reserve. In the example of Section~\ref{sec:benchmark},
the reserve lies below both class values of the rival, so every prepared rival is admissible.

A reserve that excludes the low class can widen the class difference in target proceeds, because the low class then
adds nothing to proceeds. Informed trade can then pay against both incumbents, and the price is informative at both
strengths. Proposition~\ref{prop:reserves} gives the consequence. If informed trade is full at both strengths and only a purchase recruits at each,
entry is flat in incumbent strength, and strict deterrence needs the stronger incumbent to push gross profit at the
top posterior below the high cost.

\begin{table}[tbp]
\centering\small\singlespacing
\caption{Reserves, preparation and sale in the example.}\label{tab:reserves}
\begin{tabular}{@{}llcccc@{}}
\toprule
Reserve & Strength & Entry & Sale & Two admissible bidders & Expected proceeds\\
\midrule
\(\reg{prim_p}\) & \(r_0\) & \reg{p2_reserve_points_base_r0_entry} & \reg{p2_reserve_points_base_r0_sale} & \reg{p2_reserve_points_base_r0_two} & \reg{p2_reserve_points_base_r0_proceeds}\\
 & \(r_1\) & \reg{p2_reserve_points_base_r1_entry} & \reg{p2_reserve_points_base_r1_sale} & \reg{p2_reserve_points_base_r1_two} & \reg{p2_reserve_points_base_r1_proceeds}\\
\(\reg{rsv_points_excl}\) & \(r_0\) & \reg{p2_reserve_points_exclude_r0_entry} & \reg{p2_reserve_points_exclude_r0_sale} & \reg{p2_reserve_points_exclude_r0_two} & \reg{p2_reserve_points_exclude_r0_proceeds}\\
 & \(r_1\) & \reg{p2_reserve_points_exclude_r1_entry} & \reg{p2_reserve_points_exclude_r1_sale} & \reg{p2_reserve_points_exclude_r1_two} & \reg{p2_reserve_points_exclude_r1_proceeds}\\
\(\reg{rsv_deter}\) & \(r_0\) & \reg{p2_reserve_points_deter_r0_entry} & \reg{p2_reserve_points_deter_r0_sale} & \reg{p2_reserve_points_deter_r0_two} & \reg{p2_reserve_points_deter_r0_proceeds}\\
 & \(r_1\) & \reg{p2_reserve_points_deter_r1_entry} & \reg{p2_reserve_points_deter_r1_sale} & \reg{p2_reserve_points_deter_r1_two} & \reg{p2_reserve_points_deter_r1_proceeds}\\
\bottomrule
\end{tabular}
\par\medskip
\begin{minipage}{0.92\linewidth}\footnotesize
\textit{Notes.} The example, with point-mass classes. At each reserve, full informed orders are the unique outcome
at both strengths, except at the reserve \(\reg{prim_p}\) and \(r_0\), where the informed trader does not trade.
\end{minipage}
\end{table}

Table~\ref{tab:reserves} reports three reserves in the example. The reserve
\(\reg{rsv_points_excl}\) excludes the low class. Full orders are unique at both strengths, only a purchase recruits,
and entry is \(\reg{p2_reserve_points_exclude_r0_entry}\) at both \(r_0\) and \(r_1\), while this reserve raises
expected proceeds at both strengths. Narrow uniform classes give the same pattern at the reserve
\(\reg{rsv_bands_excl}\).

The reserve \(\reg{rsv_deter}\) restores strict deterrence: entry is \(\reg{p2_reserve_points_deter_r0_entry}\) at
\(r_0\) and \(\reg{p2_reserve_points_deter_r1_entry}\) at \(r_1\), and proceeds exceed their levels at the reserve
\(\reg{prim_p}\). This reserve excludes the low class and also the weak incumbent, since it exceeds \(r_0\). In the
example, strict deterrence with a reserve that excludes the low class requires the reserve to exclude the weak
incumbent as well. With a reserve in
\((\ell,r_0]\), the top posterior recruits at both strengths, and entry is flat. Strict deterrence holds for reserves
between about \(\reg{p2_reserve_ph}\) and \(\reg{p2_reserve_weak_top}\)
(Appendix~\ref{app:welfare}).
